% filepath: c:\REPO\studia\ZAAWANSOWANAOPTYMALIZACJA20252026\ProjektSieczne\main.tex
\documentclass[10pt,a4paper]{article}

% ...existing code...
% Kodowanie i język
\usepackage[utf8]{inputenc}
\usepackage[T1]{fontenc}
\usepackage[polish]{babel}

% Układ strony i mikrotypografia
\usepackage[a4paper,margin=15mm]{geometry}
\usepackage{lmodern}
\usepackage{microtype}

% Podstawowe pakiety
\usepackage{graphicx}
\usepackage{amsmath,amssymb}
\usepackage{siunitx}
\usepackage{booktabs}
\usepackage{caption}
\usepackage{hyperref}
\usepackage{algorithm}
\usepackage{algorithmic}
% \usepackage{algpseudocode}
\usepackage{float}
\usepackage{longtable}
\usepackage{booktabs}
\usepackage{minted}
\usepackage{listings}
\usepackage{xcolor}
\usepackage{amsthm}
\newtheorem*{example}{Przykład}
\lstset{
  language=Matlab,
  basicstyle=\ttfamily\small,
  keywordstyle=\color{blue},
  commentstyle=\color{gray},
  frame=single,
  breaklines=true,
  numbers=left,
  numberstyle=\tiny,
  xleftmargin=1em
}


\hypersetup{colorlinks=true,linkcolor=blue,citecolor=blue,urlcolor=blue}

% Bibliografia (użyj biber lub zmień na natbib jeśli wolisz)
\usepackage[backend=biber,style=numeric,sorting=none]{biblatex}
\addbibresource{references.bib} % utwórz plik references.bib w katalogu projektu

% ...existing code...
\title{Norma spektralna macierzy}
\author{Paweł Drzyzga i Maciej Pasternak\\ Wydział Informatyki i Matematyki \\ Politechnika Krakowska}
\date{\today}

\begin{document}
\maketitle

\section{Norma spektralna macierzy}

Niech $\mathbb{K}\in\{\mathbb{R},\mathbb{C}\}$ oraz
\[
A\in \mathbb{K}^{m\times n},
\]
rozumiana jako operator liniowy
\[
A:(\mathbb{K}^n,\|\cdot\|_2)\longrightarrow(\mathbb{K}^m,\|\cdot\|_2),
\]
gdzie $\|\cdot\|_2$ oznacza normę euklidesową indukowaną przez standardowy iloczyn skalarny.

\subsection{Definicja}

\emph{Normą spektralną} macierzy $A$ nazywamy normę operatorową indukowaną przez normę euklidesową:
\[
\|A\|_2
=
\sup_{\|x\|_2=1}\|Ax\|_2.
\]
Jest to najmniejsza stała $C\ge 0$ taka, że
\[
\|Ax\|_2\le C\|x\|_2
\quad
\text{dla wszystkich }x\in\mathbb{K}^n.
\]

\subsection{Związek z wartościami singularnymi}

Norma spektralna jest równa największej wartości singularnej macierzy $A$:
\[
\|A\|_2=\sigma_{\max}(A).
\]
Ponieważ wartości singularne są pierwiastkami z wartości własnych macierzy $A^*A$, otrzymujemy
\[
\|A\|_2=\sqrt{\lambda_{\max}(A^*A)}.
\]

\subsection{Własności}

\subsubsection{Związek z wartościami własnymi}

Jeżeli $A$ jest operatorem normalnym (tzn.\ $AA^*=A^*A$), w szczególności gdy jest symetryczna (w $\mathbb{R}$) lub hermitowska (w $\mathbb{C}$), to
\[
\|A\|_2=\max_i|\lambda_i(A)|.
\]
W ogólności zachodzi nierówność
\[
\rho(A)=\max_i|\lambda_i(A)|\le\|A\|_2,
\]
gdzie $\rho(A)$ oznacza promień spektralny.

\subsubsection{Niezmienniczość unitarna}

Dla dowolnych macierzy unitarnych
\[
U\in\mathbb{K}^{m\times m},
\qquad
V\in\mathbb{K}^{n\times n}
\]
zachodzi
\[
\|UAV\|_2=\|A\|_2.
\]
Oznacza to, że norma spektralna zależy wyłącznie od geometrycznego działania operatora, a nie od wyboru baz w dziedzinie i przeciwdziedzinie.

\subsubsection{Podmultiplikatywność}

Dla dowolnych macierzy $A\in\mathbb{K}^{m\times n}$ i $B\in\mathbb{K}^{n\times k}$ zachodzi
\[
\|AB\|_2\le\|A\|_2\,\|B\|_2.
\]
W szczególności dla iteracji $x_{k+1}=Ax_k$ mamy
\[
\|x_k\|_2\le\|A\|_2^k\|x_0\|_2.
\]

\subsubsection{Liczba uwarunkowania}

Jeżeli $A$ jest odwracalna, to jej liczba uwarunkowania w normie spektralnej dana jest wzorem
\[
\kappa_2(A)=\|A\|_2\|A^{-1}\|_2
=\frac{\sigma_{\max}(A)}{\sigma_{\min}(A)}.
\]
Wielkość ta mierzy wrażliwość rozwiązania układu liniowego $Ax=b$ na zaburzenia danych.

\subsubsection{Relacja z normą Frobeniusa}

Dla każdej macierzy zachodzi nierówność
\[
\|A\|_2\le\|A\|_F,
\qquad
\|A\|_F=\sqrt{\sum_{i,j}|a_{ij}|^2}.
\]
Norma Frobeniusa mierzy całkowitą „energię” macierzy, natomiast norma spektralna mierzy maksymalne wzmocnienie w jednym kierunku.

\subsection{Przykłady rachunkowe}

\begin{example}[Obliczanie normy spektralnej]
Rozważmy macierz
\[
A=\begin{pmatrix}
3 & 1\\
0 & 2
\end{pmatrix}\in\mathbb{R}^{2\times 2}.
\]
Obliczymy jej normę spektralną $\|A\|_2$.

\medskip
\noindent
\textbf{Krok 1.} Obliczamy macierz $A^\top A$:
\[
A^\top=\begin{pmatrix}
3 & 0\\
1 & 2
\end{pmatrix},
\qquad
A^\top A=
\begin{pmatrix}
3 & 0\\
1 & 2
\end{pmatrix}
\begin{pmatrix}
3 & 1\\
0 & 2
\end{pmatrix}
=
\begin{pmatrix}
9 & 3\\
3 & 5
\end{pmatrix}.
\]

\medskip
\noindent
\textbf{Krok 2.} Wyznaczamy wartości własne macierzy $A^\top A$. Liczymy wielomian charakterystyczny:
\[
\det(A^\top A-\lambda I)=
\det\begin{pmatrix}
9-\lambda & 3\\
3 & 5-\lambda
\end{pmatrix}
=(9-\lambda)(5-\lambda)-9
=\lambda^2-14\lambda+36.
\]
Rozwiązując równanie
\[
\lambda^2-14\lambda+36=0,
\]
otrzymujemy
\[
\lambda=\frac{14\pm\sqrt{196-144}}{2}
=\frac{14\pm 2\sqrt{13}}{2}
=7\pm\sqrt{13}.
\]
Zatem
\[
\lambda_{\max}(A^\top A)=7+\sqrt{13}.
\]

\medskip
\noindent
\textbf{Krok 3.} Korzystając z faktu, że
\[
\|A\|_2=\sqrt{\lambda_{\max}(A^\top A)},
\]
dostajemy
\[
\|A\|_2=\sqrt{7+\sqrt{13}}.
\]

\medskip
\noindent
Jest to dokładna wartość normy spektralnej macierzy $A$.
\end{example}

\begin{example}[Macierz symetryczna: norma spektralna a wartości własne]
Rozważmy macierz symetryczną
\[
C=\begin{pmatrix}
-2 & 1\\
1 & -2
\end{pmatrix}\in\mathbb{R}^{2\times 2}.
\]
Ponieważ $C$ jest symetryczna (a więc normalna), zachodzi
\[
\|C\|_2=\max_i |\lambda_i(C)|.
\]

\medskip
\noindent
\textbf{Krok 1.} Wyznaczamy wartości własne macierzy $C$:
\[
\det(C-\lambda I)
=
\det\begin{pmatrix}
-2-\lambda & 1\\
1 & -2-\lambda
\end{pmatrix}
=
(-2-\lambda)^2-1
=
\lambda^2+4\lambda+3.
\]
Rozwiązując równanie
\[
\lambda^2+4\lambda+3=0,
\]
otrzymujemy
\[
\lambda=\frac{-4\pm\sqrt{16-12}}{2}
=\frac{-4\pm 2}{2}
\in\{-1,-3\}.
\]

\medskip
\noindent
\textbf{Krok 2.} Zatem norma spektralna wynosi
\[
\|C\|_2=\max\{|-1|,|-3|\}=3.
\]

\medskip
\end{example}


\end{document}